% !TEX root = ../main.tex
\section{Conclusiones}

El desarrollo de este laboratorio nos permitió comprobar de manera práctica cómo los principios algebraicos de la teoría de números se traducen directamente en algoritmos eficientes y robustos. Al implementar la biblioteca, confirmamos que el costo computacional de estas operaciones debe medirse siempre en función de la longitud en bits de la entrada y no de su valor numérico. En el algoritmo de Euclides, observamos que el número medio de divisiones creció a un ritmo constante de $0.584$ divisiones por bit, manteniéndose muy por debajo de la cota teórica de Lamé en la totalidad de los 7\,000 pares evaluados.

Uno de los aprendizajes más valiosos del equipo fue aprender a distinguir la cantidad de iteraciones teóricas respecto al tiempo real de ejecución medido en Python. Mientras que la cantidad de operaciones aritméticas cumple con exactitud las predicciones teóricas ---como los $k - 1$ cuadrados más los bits en 1 en la exponenciación binaria---, el tiempo observado refleja el costo intrínseco de manipular enteros de precisión arbitraria. En este sentido, la reducción modular inmediata resultó determinante: evitó la creación de cantidades intermedias inmanejables de millones de bits y mantuvo acotada la memoria durante todo el cómputo.

Asimismo, verificar el Teorema Chino del Resto y las congruencias lineales nos demostró que controlar las condiciones algebraicas de existencia, como la coprimalidad y el valor del MCD, es tan crucial como la propia fórmula matemática de cálculo. Comprobar que más de 12\,000 casos aleatorios coincidieron plenamente con las funciones nativas y con las soluciones esperadas nos da la seguridad de contar con una base numérica confiable y lista para su aplicación en criptografía y algoritmos avanzados.
